% Part: methods
% Chapter: proofs
% Section: example-2

\documentclass[../../../include/open-logic-section]{subfiles}

\begin{document}

\olfileid[id]{mth}{prf}{ex2}

\olsection{Contoh Lain}

\begin{prop}
Jika $A \subseteq C$, maka $A \cup (C \setminus A) = C$.
\end{prop}

\begin{proof}
  Andaikan $A \subseteq C$. Kita ingin menunjukkan bahwa $A \cup (C
  \setminus A) = C$.
  \begin{quote}
    Kita mulai dengan mengamati bahwa ini merupakan pernyataan
    kondisional. Pernyataan itu secara tersirat dikuantifikasi secara
    universal: proposisi berlaku untuk semua himpunan $A$ dan $C$.
    Jadi, $A$ dan $C$ adalah variabel bagi himpunan sembarang. Untuk
    membuktikan pernyataan semacam itu, kita mengasumsikan antesedennya
    dan membuktikan konsekuennya.

    Kita melanjutkan dengan menggunakan asumsi bahwa $A \subseteq C$.
    Mari kita uraikan definisi~$\subseteq$: asumsi tersebut berarti
    bahwa semua !!{element} dari~$A$ juga merupakan !!{element}
    dari~$C$. Mari kita tuliskan hal ini---ini adalah fakta penting yang
    akan kita gunakan di sepanjang bukti.
  \end{quote}
  Menurut definisi~$\subseteq$, karena $A \subseteq C$, untuk setiap
  $z$, jika $z \in A$, maka $z \in C$.
  \begin{quote}
    Kita telah menguraikan semua definisi yang diberikan dalam asumsi.
    Sekarang kita dapat beralih ke kesimpulan. Kita ingin menunjukkan
    bahwa $A \cup (C \setminus A) = C$, sehingga kita menyiapkan bukti
    dengan cara yang serupa dengan contoh sebelumnya: kita menunjukkan
    bahwa setiap !!{element} dari $A \cup (C \setminus A)$ juga
    merupakan !!a{element} dari~$C$, dan sebaliknya, setiap !!{element}
    dari $C$ merupakan !!a{element} dari $A \cup (C \setminus A)$. Kita
    dapat menyingkatnya menjadi: $A \cup (C \setminus A) \subseteq C$
    dan $C \subseteq A \cup (C \setminus A)$. (Di sini kita melakukan
    kebalikan dari menguraikan definisi, tetapi hal itu membuat bukti
    sedikit lebih mudah dibaca.) Karena ini merupakan konjungsi, kita
    harus membuktikan kedua bagiannya. Untuk menunjukkan bagian
    pertama, yaitu bahwa setiap !!{element} dari $A \cup (C \setminus
    A)$ juga merupakan !!a{element} dari~$C$, kita mengasumsikan $z
    \in A \cup (C \setminus A)$ untuk suatu~$z$ sembarang, lalu
    menunjukkan bahwa $z \in C$. Menurut definisi $\cup$, kita dapat
    menyimpulkan bahwa $z \in A$ atau $z \in C \setminus A$ dari $z
    \in A \cup (C \setminus A)$. Sekarang Anda semestinya mulai
    memahami polanya.
  \end{quote}
  $A \cup (C \setminus A) = C$ jika dan hanya jika $A \cup (C
  \setminus A) \subseteq C$ dan $C \subseteq A \cup (C \setminus A)$.
  Pertama-tama kita buktikan bahwa $A \cup (C \setminus A) \subseteq
  C$. Misalkan $z \in A \cup (C \setminus A)$. Jadi, $z \in A$ atau
  $z \in (C \setminus A)$.
  \begin{quote}
    Kita telah sampai pada suatu disjungsi, dan darinya kita ingin
    membuktikan bahwa $z \in C$. Kita melakukannya dengan pembuktian
    berdasarkan kasus.
  \end{quote}
  Kasus 1: $z \in A$. Karena untuk setiap $z$, jika $z \in A$, maka $z
  \in C$, kita memperoleh $z \in C$.
  \begin{quote}
    Di sini kita menggunakan fakta yang dicatat sebelumnya, yang
    mengikuti hipotesis proposisi bahwa $A \subseteq C$. Kasus pertama
    selesai, dan kita beralih ke kasus kedua, $z \in (C \setminus A)$.
    Ingat bahwa $C \setminus A$ menyatakan \emph{selisih} kedua
    himpunan, yaitu himpunan semua !!{element} dari~$C$ yang bukan
    !!{element} dari~$A$. Namun, setiap !!{element} dari $C$ yang tidak
    berada dalam~$A$ khususnya merupakan !!a{element} dari~$C$.
  \end{quote}
  Kasus 2: $z \in (C \setminus A)$. Ini berarti bahwa $z \in C$ dan $z
  \notin A$. Jadi, khususnya, $z \in C$.
  \begin{quote}
    Bagus, kita telah membuktikan arah pertama. Sekarang kita beralih
    ke arah kedua. Di sini kita membuktikan bahwa $C \subseteq A \cup
    (C \setminus A)$. Jadi, kita mengasumsikan bahwa $z \in C$ dan
    membuktikan bahwa $z \in A \cup (C \setminus A)$.
  \end{quote}
  Sekarang misalkan $z \in C$. Kita ingin menunjukkan bahwa $z \in A$
  atau $z \in C \setminus A$.
  \begin{quote}
    Karena semua !!{element} dari $A$ juga merupakan !!{element} dari
    $C$, dan $C \setminus A$ adalah himpunan semua objek yang merupakan
    !!{element} dari $C$ tetapi bukan dari $A$, maka $z$ berada dalam
    $A$ atau dalam $C \setminus A$. Hal ini mungkin sedikit kurang
    jelas jika Anda belum mengetahui mengapa hasilnya benar. Akan
    lebih baik jika kita membuktikannya langkah demi langkah. Akan
    membantu jika kita menggunakan sebuah fakta sederhana yang dapat
    kita nyatakan tanpa bukti: $z \in A$ atau $z \notin A$. Ini disebut
    ``prinsip tengah yang dikecualikan'': untuk sebarang pernyataan~$p$,
    $p$ benar atau negasinya benar. (Di sini, $p$ adalah pernyataan
    bahwa $z \in A$.) Karena ini merupakan disjungsi, kita kembali
    dapat menggunakan pembuktian dengan kasus.
  \end{quote}
  $z \in A$ atau $z \notin A$. Dalam kasus pertama, $z \in A \cup (C
  \setminus A)$. Dalam kasus kedua, $z \in C$ dan $z \notin A$, jadi
  $z \in C \setminus A$. Namun, dengan demikian $z \in A \cup (C
  \setminus A)$.
  \begin{quote}
    Bukti kita selesai: kita telah menunjukkan bahwa $A \cup (C
    \setminus A) = C$.
  \end{quote}
\end{proof}

\end{document}
